%======================================================================
%  From the specification to LLM inference
%======================================================================

\section{What the platform implies for LLM inference}
\label{sec:platform-llm}

The preceding sections describe the platform; this one states what its
specification implies for the workload studied in the thesis.  Every
number here is derived from nominal peaks (\tabref{tab:xeon-max-peaks})
and from model architecture, not measured.  The next chapter replaces the
nominal ceilings with measured ones and builds the measured Rooflines of
both memory configurations; the chapter after it compares what two
inference back ends, \lcpp{} and vLLM, actually achieve against those
ceilings.

\subsection{Where inference sits on the Roofline}
\label{sec:platform-nominal-roofline}

\tabref{tab:nominal-ridge} gives the ridge point of each compute path
against each memory tier, the arithmetic intensity at which a kernel stops
being limited by bandwidth (\secref{sec:roofline}).  A matrix-vector
product in decode performs approximately two operations per matrix
weight for each evaluated sequence.  If the matrix is fetched once
from main memory and other traffic is neglected, then with \(b\)
evaluated sequences and \(w\) bytes per weight its intensity is about
\(2b/w\): roughly \(b\) FLOP/byte for BF16 and \(2b\) operations per
byte for INT8.  Prefill of \(P\) positions evaluated together gives the
analogous weight-only estimate of \(P\) FLOP/byte in BF16~\cite{na2024llmcpu}.
These estimates describe dense products, not every parameter access or
the complete phase (\secref{sec:memory-requirements}).

\begin{table}[htbp]
\centering
\small
\tabsetup
\begin{tabularx}{\textwidth}{P{3.4cm}P{2.6cm}YY}
\headrow
\headcell{Compute path} & \headcell{Peak per socket} &
\headcell{Ridge vs HBM (1\,TB/s)} & \headcell{Ridge vs DDR5 (307.2\,GB/s)} \\[3pt]
\tabheadrule
\tabstripeRow
AVX-512 FP32 & 6.8\,TFLOP/s & 6.8\,FLOP/byte & 22\,FLOP/byte \\[3pt]

\tabrowsep
\tabplainRow
AMX BF16 & 108.9\,TFLOP/s & 109\,FLOP/byte & 355\,FLOP/byte \\[3pt]

\tabrowsep
\tabstripeRow
AMX INT8 & 217.9\,TOP/s & 218\,OP/byte & 709\,OP/byte \\[3pt]
\tabbottomrule
\end{tabularx}
\caption{Nominal ridge points of one Max 9480 socket, from the peaks in
\tabref{tab:xeon-max-peaks}.  Integer paths use operations rather than
FLOP.  Measured ridge points follow in the next chapter.}
\label{tab:nominal-ridge}
\end{table}

Under the weight-only assumptions, single-sequence BF16 decode has an
intensity of about one FLOP per byte, two orders of magnitude left of
the nominal AMX ridge on either tier; even 32 evaluated sequences remain
left of it on HBM.  This suggests a bandwidth constraint for these dense
products if the assumed main-memory traffic occurs, rather than proving
the bottleneck of every serving configuration.  Attention additionally
accesses KV histories, generally distinct across sequences, with reuse
depending on heads, sharing and tiling.  Prefill of a few hundred positions
evaluated together can approach or cross the nominal AMX ridge in this
estimate.  Actual phase intensity requires the measured traffic and
kernel execution path.

\subsection{What fits where}
\label{sec:platform-capacity}

Capacity decides which tier can hold a model at all.
\tabref{tab:model-capacity} gives the weight footprint of two
representative model sizes with the Llama~3.1 architecture, 8 and 70
billion parameters~\cite{grattafiori2024llama3}, at three precisions,
against the HBM available at the three levels of locality: one tile under
SNC4 (16\,GB), one socket (64\,GB) and one node (128\,GB).

\begin{table}[htbp]
\centering
\small
\tabsetup
\begin{tabularx}{\textwidth}{P{2.2cm}P{2.6cm}P{2.0cm}Y}
\headrow
\headcell{Model} & \headcell{Weight precision} & \headcell{Weights} &
\headcell{Smallest HBM domain that holds them} \\[3pt]
\tabheadrule
\tabstripeRow
8B & BF16 & 16.1\,GB & One tile only nominally (15.0 of 16\,GiB); one
  socket once KV cache and buffers are added \\[3pt]

\tabrowsep
\tabplainRow
8B & INT8 & 8.0\,GB & One tile \\[3pt]

\tabrowsep
\tabstripeRow
8B & 4-bit (4.8\,bits/weight) & 4.9\,GB & One tile \\[3pt]

\tabrowsep
\tabplainRow
70B & BF16 & 141\,GB & None; exceeds the node's 128\,GB \\[3pt]

\tabrowsep
\tabstripeRow
70B & INT8 & 70.6\,GB & Both sockets, split between them \\[3pt]

\tabrowsep
\tabplainRow
70B & 4-bit (4.8\,bits/weight) & 42.5\,GB & One socket \\[3pt]
\tabbottomrule
\end{tabularx}
\caption{Weight footprints of representative model sizes against HBM
capacity.  The 4-bit row uses the effective size of common 4-bit GGUF
formats, including their scale factors~\cite{llamacpp2026}.  Usable
capacity is lower than nominal: in HBM-only mode the operating system
shares it, and KV cache and runtime buffers compete for it in every
mode.}
\label{tab:model-capacity}
\end{table}

The KV cache adds a second, workload-dependent footprint.  With grouped-query
attention, eight KV heads of dimension 128 and BF16 entries, one token
occupies \(2 \times 32 \times 8 \times 128 \times 2\)\,bytes, or 128\,KiB,
in the 8B model's 32 layers, and 320\,KiB in the 70B model's 80 layers.  A
single 32k-token context therefore needs 4\,GiB and 10\,GiB respectively,
and sixteen concurrent 8k-token contexts need 16\,GiB and 40\,GiB.  An 8B
model in BF16 plus such a batch already fills half of a socket's HBM; a
70B model in 4-bit precision plus the same batch does not fit in it.  This
is the regime in which deciding what goes in HBM, weights or KV, becomes
a real choice rather than a formality.

\subsection{Upper bounds on decode throughput}
\label{sec:platform-token-bounds}

For single-sequence decode, let \(D_W\) and \(D_{\mathrm{KV}}\) be
main-memory bytes read per evaluated position for weights and KV state.
If both use the same memory boundary with bandwidth ceiling \(B\),
neglecting other traffic gives the conditional bound

\begin{equation}
\label{eq:token-bound}
  r_{\text{token}} \leq \frac{B}{D_W + D_{\mathrm{KV}}}.
\end{equation}
\eqcaption{Conditional bandwidth bound on single-sequence decode:
bandwidth divided by weight and KV read traffic per evaluated position.
Payload size is not automatically equal to this traffic.}

\tabref{tab:token-bounds} additionally assumes \(D_W\approx M_W\), a
full weight representation fetch per position, and neglects KV traffic
for a short context.  Cache reuse, embedding lookups and representation
details can invalidate this approximation.  Under these assumptions,
replacing the DDR5 ceiling with the HBM ceiling changes the estimate by
their nominal ratio, about 3.3; this is not a demonstrated placement gain.
Data split across tiers requires separate traffic and path accounting.
Sustainable bandwidth is lower than the interface peaks, and the runtime
may achieve only part of it; measurements establish that utilization.

\begin{table}[htbp]
\centering
\small
\tabsetup
\begin{tabularx}{\textwidth}{P{3.6cm}P{2.2cm}YY}
\headrow
\headcell{Model and precision} & \headcell{Weights} &
\headcell{Ceiling from HBM (1\,TB/s)} &
\headcell{Ceiling from DDR5 (307.2\,GB/s)} \\[3pt]
\tabheadrule
\tabstripeRow
8B, BF16 & 16.1\,GB & 62\,tokens/s & 19\,tokens/s \\[3pt]

\tabrowsep
\tabplainRow
8B, 4-bit & 4.9\,GB & 203\,tokens/s & 62\,tokens/s \\[3pt]

\tabrowsep
\tabstripeRow
70B, 4-bit & 42.5\,GB & 24\,tokens/s & 7.2\,tokens/s \\[3pt]

\tabrowsep
\tabplainRow
70B, INT8 & 70.6\,GB & Does not fit one socket & 4.4\,tokens/s \\[3pt]
\tabbottomrule
\end{tabularx}
\caption{Nominal single-socket upper bounds on the decode rate of one
sequence, from \eqnref{eq:token-bound} assuming one full weight
representation fetch per position and neglecting KV and other traffic.
Splitting a model across both sockets can at most double a bound, and
only if each socket reads its half locally.}
\label{tab:token-bounds}
\end{table}

\subsection{Concurrency needed to reach the bandwidth}
\label{sec:platform-concurrency}

Little's law (\eqnref{eq:little}) converts these bandwidths into a
demand on the cores.  The latency to use is the loaded latency of the
platform, which the next chapter measures.  An optimistic lower bound
follows from the idle latencies Ibeid et al.\ measured on a Xeon Max
9470C, 134.9\,ns for HBM and 112.9\,ns for DDR5~\cite{ibeid2025aurora}.
Sustaining the nominal 1\,TB/s of HBM against 134.9\,ns requires about
135\,KB in flight, some 2{,}100 cache lines, or 38 per core across 56
cores.  Sustaining DDR5's 307.2\,GB/s against 112.9\,ns requires about
540 lines, or 10 per core.  Under load the latencies grow and so do both
requirements.

Thirty-eight outstanding lines per core is more than twice the 16 L1 fill
buffers of a Golden Cove core and close to the 48 misses it can track in
total (\secref{sec:xeon-max-cores}).  Approaching HBM bandwidth therefore
needs nearly every core of the socket, with the L2 hardware prefetchers
generating most of the requests, and an access pattern regular enough for
them to predict.  The weight stream of decode is such a pattern; the
dependent chain through the layers and the attention reads over a
growing KV cache are less so.  Two consequences follow for the
experiments.  The number of cores is a first-class factor, not a
constant set to ``all cores''.  And a disappointing HBM result has to be
explained by the bandwidth actually achieved before it is reported as an
absence of effect.
